\documentclass[11pt]{article}
\usepackage[margin=1in]{geometry}
\usepackage{amsmath,amssymb,amsthm,booktabs,hyperref}
\newtheorem{theorem}{Theorem}
\newtheorem{lemma}{Lemma}
\newtheorem{proposition}{Proposition}
\title{Finite-template refinements for the Sidorenko construction}
\author{AI-generated research draft prepared with Codex}
\date{October 10, 2026}
\begin{document}
\maketitle
\paragraph{Scope and attribution.}
This is an independent follow-up to
\href{https://github.com/openai/math/tree/fd4aeeb2ee4fc729c18d98444fed42fd0529eeeb/preprints/A-counterexample-to-Sidorenkos-conjecture-September-23-2026}{OpenAI's original manuscript},
pinned at commit \texttt{fd4aeeb2ee4fc729c18d98444fed42fd0529eeeb}.
It is not an official OpenAI publication. Mathematical priority and
independent formal verification are not claimed. Exact checks supplement
the proofs below and do not verify the imported source lemmas.
\section{Sidorenko: a structural reduction and a certified finite minimum}
\label{sec:sidorenko-minimum}
This section refers to OpenAI's
\href{https://github.com/openai/math/tree/fd4aeeb2ee4fc729c18d98444fed42fd0529eeeb/preprints/A-counterexample-to-Sidorenkos-conjecture-September-23-2026}{Sidorenko manuscript},
particularly its finite-complex proposition and the exposure and pinning arguments.
We now allow any finite probability distribution on admissible cuts,
not just three equally weighted cuts. All faces remain distinct triples.
An admissible cut partitions the faces into two classes, each exposing
all $n$ points by starting with one triangle and successively adding a
triangle sharing a pair with the existing complex and introducing one
new point. Every point pair occurring in the full complex belongs to
exactly two faces.

\begin{lemma}\label{lem:hamilton-reduction}
The existence of one admissible cut makes the complex a simple
triangulation of the sphere with $2n-4$ faces and $3n-6$ point pairs.
The point pairs crossed by each admissible cut form a Hamiltonian cycle
in its point graph. Conversely every Hamiltonian cycle in this sphere
triangulation gives an admissible cut.
\end{lemma}
\begin{proof}
In one exposure class, the shared edge of an added face is on the boundary:
an interior edge already has its two incident faces and cannot belong to
another face. Its third vertex is new. Starting from a triangle, this
operation preserves a triangulated disk whose boundary contains every
vertex exactly once. The class has $n-2$ triangles and a boundary cycle
of length $n$. An edge is internal to this disk if and only if its two
incident faces are in this class. The two classes therefore have exactly
the same boundary edges and no shared interior edges. Gluing the disks
along their common boundary gives a triangulated sphere. Distinct triples
and pairs give a simple point graph; the usual sphere counts follow.
The common boundary is precisely the cut's crossing set.
Conversely a Hamiltonian cycle separates the sphere into two triangulated
polygons without interior vertices. The dual graph of each polygon's
triangulation is a tree. Removing leaf triangles and then reversing that
order gives the required exposure. Each removed leaf has a vertex which
occurs in no remaining triangle, so each reverse step introduces one
new point.
\end{proof}

Let $a_e$ be the probability that a random admissible cut crosses $e$.
Every boundary cycle has degree two at each point, whence
\[
 \sum_{e\ni v}a_e=2,\qquad \sum_ea_e=n.
\]
The source's requirement $a_e\ge1/3$ consequently implies
$\deg(v)\le6$. At a degree-six point every incident edge must have
$a_e=1/3$. Also
\[
 \sum_e(a_e-1/3)=n-(3n-6)/3=2,
\]
so strict inequality somewhere follows automatically.

\begin{proposition}[Computer-assisted minimum within the finite template]
Subject to completeness of the plantri enumeration and the exact checks
described below, no distinct-triple complex on fewer than $13$ points
simultaneously has the paired exposure property, two occurrences of every
point pair, a distribution of cuts with $a_e\ge1/3$, and discrete color
refinement of its incidence graph starting from its two part colors.
Thus the smallest incidence graph in this template has $35$ vertices.
The original $13$-point, $22$-face certificate attains that size.
\end{proposition}
\begin{proof}[Enumeration and certificate argument]
By Lemma~\ref{lem:hamilton-reduction}, it suffices to enumerate simple
sphere triangulations, first filter by maximum degree six, and then
apply incidence color refinement. For $n\le3$ there cannot be the
required distinct $2n-4$ faces with every pair in two faces. For
$4\le n\le12$, plantri 5.8 gives the following counts:
\begin{center}\small
\begin{tabular}{@{}rrrrr@{}}
\toprule
$n$&All triangulations&Degree at most six&Discrete refinement&Obstructions\\
\midrule
4&1&1&0&0\\
5&1&1&0&0\\
6&2&2&0&0\\
7&5&5&0&0\\
8&14&10&0&0\\
9&50&15&1&1\\
10&233&30&4&4\\
11&1249&44&12&12\\
12&7595&77&23&23\\
\bottomrule
\end{tabular}
\end{center}
For each of the $40$ surviving graphs, a rational vector $y_e\ge0$,
$\sum_ey_e=1$, satisfies
\[
 \max_{C\text{ Hamiltonian}}\sum_{e\in C}y_e<1/3.
\]
A putative cut distribution would instead give
$\mathbb E_C\sum_{e\in C}y_e=\sum_ey_ea_e\ge1/3$, a contradiction.
The smallest verified separation margin is $1/75$.

The proposal computation enumerates cycles and uses floating-point linear
programming only to suggest edge weights. A separate verifier uses exact
rational arithmetic and Held--Karp dynamic programming to maximize cycle
weight and count all Hamiltonian cycles, without importing the proposal
or exposure code. It checks $16{,}941$ cycles in aggregate, reconstructs
faces independently from rotations, repeats refinement by partition
splitting, and rejects a corrupted witness. Generator output counts and
hashes are retained. This is a computer-assisted certificate conditional
on the generator's completeness, not a Lean proof or a minimality theorem
for arbitrary Sidorenko counterexamples.
\end{proof}

The enumeration also inspected all $49{,}566$ triangulations at $n=13$.
Of $115$ with degree at most six, $49$ have discrete incidence refinement.
The cut feasibility pass produced $47$ exact obstruction witnesses and
two exactly checked feasible finite objects. One is isomorphic to the
original complex; the other is not, as distinguished by the canonical
adjacency encoding from their discrete refinement colors. Both have
$35$ incidence vertices and $66$ edges. The alternative has not undergone
an independent transfer of the full analytic kernel argument and is not
claimed here as a newly proved Sidorenko counterexample.

\section{Sidorenko: a smaller nondirect residual bound}
\label{sec:residual-improvement}
The original tail section's lemma ``Costs for the remaining profiles''
bounds the nondirect contribution by $O_D(q^{-D+1641})$. Its equality
profiles already require the excess-separation term
\[
 \mathcal E=\tfrac12\sum_e(a_e-1/3)b(h_e)
\]
to vanish, with every $h_e$ either zero or two. Hence $h_e=0$ whenever
$a_e>1/3$. For the original manuscript's exact three bit cuts, those
five excluded pairs are
\[
 \{0,2\},\ \{7,12\},\ \{4,8\},\ \{5,6\},\ \{5,10\}.
\]
Only $28$ pairs can therefore be active. The degrees in this allowed
pair graph, in the original point order, are
\[
 (3,6,3,6,4,2,4,4,4,6,4,6,4).
\]
This uses the published cuts, without changing the graph or its laws.

\begin{proposition}\label{prop:residual-1032}
Within the source's nondirect residual stratum, for fixed even
$D\ge1032$, the contribution is $O_D(q^{-D+1032})$.
In particular that stratum vanishes for even $D\ge1034$.
All other dimension thresholds in the source remain required.
\end{proposition}
\begin{proof}
Retain the source's notation: $m$ active pairs, global span dimension $k$,
local span dimensions $k_i$, and defect
$d=\sum_i k_i-2k\ge1$. Including the source's gain bound $2m$, its
counting exponent is
\[
 -Dd+R,\qquad
 R=-k^2+2mk-2m+\sum_i\binom{k_i}{2}.
\]
Nondirectness implies $2\le k\le2m-1$, and now $m\le28$.
The local caps are twice the displayed degrees, so
\[
 \sum_i\binom{k_i}{2}\le468,\qquad
 R\le m^2-2m+468\le1196.
\]
For $d\ge2$ this is already sufficient, since for $D\ge1032$,
$-Dd+R\le-2D+1196\le-D+1032$.

It remains to handle $d=1$, when $\sum_i k_i=2k+1$.
Sort the local caps into
\[
 c=(12,12,12,12,8,8,8,8,8,8,6,6,4).
\]
For total $S$, let $G(S)$ be the sum of $\binom{x_i}{2}$ obtained by
filling these caps in order until total $S$ is reached. This bounds every
feasible local sum: the sorted greedy vector majorizes any vector with
these caps and total, and $x\mapsto\binom{x}{2}$ is convex.
Since $2k-2>0$, replacing $m$ by $28$ only increases the remainder.
Consequently
\[
 R\le-k^2+56k-56+G(2k+1)\le1032
 \quad(2\le k\le55).
\]
The last inequality is a finite integer calculation over $54$ cases;
its maxima occur at $k=31,32$. For example $G(63)=313$ and $G(65)=320$,
giving $1032$ in both cases. The exact audit checks every case and
independently computes $G$ by integer dynamic programming for every
possible total. It also exhausts all relaxed positive defects, obtaining
the same maximum $R/d=1032$. Thus each stratum has the asserted exponent.
There are only finitely many strata for fixed $D$, as in the source.
\end{proof}

This improves the sufficient threshold for this one counting step from
$1642$ to $1034$ among even dimensions. It is not a sufficient dimension
for the entire rank-layer theorem, whose earlier residue and transverse
thresholds may dominate. The arithmetic certificate is supplementary to
the proof and does not verify those imported arguments.


\section{Active-support refinement of the residual bound}
Let $S$ consist of the vertices incident to an active pair space, and
let $m_S$ count allowed graph edges induced by $S$. Each active pair
space has dimension two. Therefore
\[
 m\le m_S,\qquad 2\le k_i\le\min\{k,2\deg_S(i)\}\quad(i\in S),
 \qquad k_i=0\quad(i\notin S).
\]
For defect one, $\sum_i k_i=2k+1$. Since $2k-2>0$, the remainder is
at most $-k^2+2m_Sk-2m_S+\sum_i\binom{k_i}{2}$.
Sort the caps decreasingly, start each supported coordinate at two,
then fill the largest caps first to the required total. This vector
majorizes every feasible vector: the first $j$ sorted entries are bounded
by both the first-$j$ cap sum and the total minus two for each remaining
entry. The greedy vector attains these bounds. Convexity gives the
maximum binomial sum. An independent integer convolution confirms it.

The support audit checks all $8192$ vertex subsets and $55{,}173$
feasible relaxed defect-one cases, obtaining a maximum remainder of
$978$. A maximizing relaxation has full support, $m_S=28$, $k=33$,
and dimensions $(12,12,12,12,3,2,2,2,2,2,2,2,2)$; its binomial sum
is $275$, so $-33^2+56\cdot33-56+275=978$.
Realizability of the maximizing relaxation is not required for an upper
bound. For defect at least two, the earlier bound $1196$ yields
$-2D+1196\le-D+978$ when $D\ge978$. Consequently this one residual
contribution is $O_D(q^{-D+978})$ and vanishes for even $D\ge980$.
All earlier thresholds in the original analytic argument remain required.
The detailed derivation is in \texttt{support-refinement.md}; the audit
is \texttt{audit\_support.py}, now included in the standard check runner.

\section{Sequential basis count for the nondirect residual}
\label{sec:sequential}
The two preceding bounds count the $m$ active pair spaces as free
two-dimensional subspaces of the global span. The original direct-case
count is sharper: it places one basis vector at a time and uses that every
new vector is symplectically orthogonal to everything already placed at
either endpoint. That count extends to nondirect families once each basis
vector records whether it raises the global span and whether it raises the
two endpoint spans. The orthogonality conditions then cancel the
containment cross terms exactly as in the direct case, and the defect
$d=\sum_ik_i-2k$ counts the vectors that raise a local span without raising
the global one.

Fix an active pair set $E'$ with $|E'|=m$, an order $e_1,\dots,e_m$ on it,
and a family $H_t\subseteq L_i\cap L_j$ of two-dimensional subspaces,
$e_t=\{i,j\}$. Choose an ordered basis $(u_t,w_t)$ of each $H_t$ and list
the vectors as $u_1,w_1,u_2,w_2,\dots$. For a vector $x$ of $H_t$ let
$W'(x)$ be the span of all earlier vectors in this list, and for an endpoint
$i$ of $e_t$ let $V'_i(x)$ be the span of the earlier vectors belonging to
pairs containing $i$. Write $w'(x)$, $v'_i(x)$, $v'_j(x)$ for their
dimensions and $\pi(x)=\dim(V'_i(x)\cap V'_j(x))$. The \emph{type} of $x$ is
$(g,\ell_i,\ell_j)\in\{0,1\}^3$, where $g=1$ if $x\notin W'(x)$ and
$\ell_i=1$ if $x\notin V'_i(x)$.

\begin{lemma}\label{lem:sequential-accounting}
With this notation, $\ell_i,\ell_j\ge g$ for every vector,
\[
 k=\#\{x:g(x)=1\},\qquad k_i=\#\{x\text{ at }i:\ell_i(x)=1\},\qquad
 d=\sum_x\bigl(\ell_i(x)+\ell_j(x)-2g(x)\bigr),
\]
and
\begin{equation}\label{eq:sequential-identity}
 \sum_i\binom{k_i}{2}=\sum_x\sum_{\text{endpoints }i\text{ of }x\text{ with }\ell_i(x)=1}v'_i(x).
\end{equation}
Every vector is orthogonal to $V'_i(x)+V'_j(x)$, and
$\pi(x)\le\min\{v'_i(x),v'_j(x)\}$.
\end{lemma}
\begin{proof}
Since $V'_i(x)\subseteq W'(x)$, a vector inside the earlier global span may
still be outside an earlier local span, but not conversely. Each vector
with $g=1$ raises the global span by exactly one, and each vector with
$\ell_i=1$ raises the span at $i$ by exactly one; the other vectors change
nothing. This gives the three counts, and the defect formula is their
combination. At a vertex $i$, the vectors with $\ell_i=1$ see earlier local
dimensions $0,1,\dots,k_i-1$ in order, whose sum is $\binom{k_i}{2}$;
summing over vertices gives Equation~\eqref{eq:sequential-identity}.
Orthogonality holds because $x$ and $V'_i(x)$ lie in the Lagrangian
$L_i$, and likewise at $j$. The intersection bound is trivial.
\end{proof}

\begin{proposition}\label{prop:residual-189}
Within the source's nondirect residual stratum, for fixed even $D\ge189$,
the contribution is $O_D(q^{-D+189})$. In particular that stratum vanishes
for even $D\ge190$. All other dimension thresholds in the source remain
required.
\end{proposition}
\begin{proof}
Fix a profile, the activity data, and the stratum $(k,k_i)$. Refine it by
the sequence of types, a finite set of patterns independent of $D$. Count
the vector lists of a fixed pattern by choosing the vectors in order.
A vector of type $(1,1,1)$ is orthogonal to $V'_i(x)+V'_j(x)$, a space of
dimension $v'_i+v'_j-\pi$, so by nondegeneracy there are exactly
$q^{2D-v'_i-v'_j+\pi}$ candidates. A vector of type $(0,0,0)$ lies in
$V'_i\cap V'_j$: at most $q^{\pi}$ candidates. A vector of type $(0,0,1)$
lies in $V'_i$: at most $q^{v'_i}$ candidates, and symmetrically for
$(0,1,0)$. A vector of type $(0,1,1)$ lies in $W'$: at most $q^{w'}\le q^k$
candidates. Each pair space has at least $\tfrac38q^4$ ordered bases, so
the number of families with this pattern is at most $(8/3)^mq^{-4m}$
times the product of these counts. Multiplying by the containment
probabilities $O_D(q^{-Dk_i+\binom{k_i}{2}})$ for the isotropic local spans,
as in the source, and by the gain $q^{2m}$, the exponent is
$2Dk-D\sum_ik_i$ plus
\[
 R=\sum_i\binom{k_i}{2}-2m
 -\sum_{(1,1,1)}(v'_i+v'_j-\pi)
 +\sum_{(0,0,0)}\pi+\sum_{(0,0,1)}v'_i+\sum_{(0,1,1)}w',
\]
where the $(0,1,0)$ terms are included in the third sum by symmetry.
Equation~\eqref{eq:sequential-identity} evaluates
$\sum_{(1,1,1)}(v'_i+v'_j)$ as $\sum_i\binom{k_i}{2}$ minus
$\sum_{(0,0,1)}v'_j$ minus $\sum_{(0,1,1)}(v'_i+v'_j)$. Hence
\begin{equation}\label{eq:sequential-remainder}
 R=-2m+\sum_{(1,1,1),(0,0,0)}\pi
 +\sum_{(0,0,1)}(v'_i+v'_j)+\sum_{(0,1,1)}(w'+v'_i+v'_j).
\end{equation}
The defect is the number of $(0,0,1)$ vectors plus twice the number of
$(0,1,1)$ vectors, so the exponent is $-Dd+R$ and the first two sums run
over all but $d$ vectors at most.

Bound $\pi\le\tfrac12(v'_i+v'_j)$ on every vector. At a vertex $i$ of
active degree $\delta_i$, the $k_i$ vectors with $\ell_i=1$ contribute
$\binom{k_i}{2}$ to $\sum_xv'_i(x)$, and the remaining $2\delta_i-k_i$
vectors contribute at most $k_i$ each. Therefore
\[
 \sum_x\bigl(v'_i(x)+v'_j(x)\bigr)
 \le\sum_i\varphi_i(k_i),\qquad
 \varphi_i(x)=\binom{x}{2}+(2\delta_i-x)x
 \le\delta_i(2\delta_i-1),
\]
the last inequality by concavity, with equality exactly at
$x\in\{2\delta_i-1,2\delta_i\}$. For $d=1$ the single defect vector has
type $(0,0,1)$ and its endpoints see $v'_i\le k_i$ and $v'_j\le k_j-1$. So
\[
 R\le-2m+\frac12\sum_i\delta_i(2\delta_i-1)+\frac{k_i+k_j-1}{2}
 =\sum_i\delta_i^2-3m+\frac{k_i+k_j-1}{2}.
\]
Adding a pair whose endpoints currently have degrees $a,b$ changes
$\sum_i\delta_i^2-3m$ by $2a+2b-1$, which is negative only when both
endpoints are isolated. The allowed graph is connected, so the full
allowed graph maximizes this quantity: with degrees
$(3,6,3,6,4,2,4,4,4,6,4,6,4)$ and $m=28$ it equals $262-84=178$. The
defect pair joins two points of degree at most six, so
$k_i+k_j-1\le23$. As $R$ is an integer, $R\le178+11=189$.

For $d\ge2$, write $d=n_1+2n_2$ with $n_1$ vectors of type $(0,0,1)$ and
$n_2$ of type $(0,1,1)$. The same averaging gives
$R\le178+\tfrac{23}{2}n_1+(k+11)n_2$ with $k\le2m-1=55$, so $R/d\le122$.
Thus $-Dd+R\le-D+189$ for every $d\ge1$ once $D\ge189$. There are finitely
many strata and patterns at fixed $D$, as in the source.
\end{proof}

The exact audit \texttt{audit\_sequential.py} reconstructs the allowed
graph from the saved faces and cut bits, checks its connectivity and the
per-vertex concavity bound, and exhausts all $8192$ supports and
$55{,}173$ relaxed defect-one strata, computing the local maximum by an
integer convolution and independently by a marginal-gain allocation. Its
maximum remainder is $189$, attained only on the full support with every
local dimension at or one below full rank, matching the closed form. It
also verifies the defect-at-least-two ratio $122$. Even dimension $190$
therefore suffices for this one counting step, instead of $980$.

\section{Overlaps come only from defect vectors}
\label{sec:overlap}
The averaging step at the end of Section~\ref{sec:sequential} bounds the
overlap $\pi(x)=\dim(V'_i(x)\cap V'_j(x))$ by the smaller endpoint span.
The overlap is in fact controlled by the defect alone. For a vector $y$ at a
pair containing the point $a$ write $\tau_a(y)=\ell_a(y)-g(y)\in\{0,1\}$,
so that $\varepsilon(y)=\tau_i(y)+\tau_j(y)$ is its contribution to $d$.

\begin{lemma}\label{lem:overlap}
Let $x$ be a basis vector of the pair $e_t=\{i,j\}$ and let $F$ be the set
of earlier pairs sharing a point with $e_t$. Then
\[
 \pi(x)\le\sum_{y\in F}\tau_{a(y)}(y)
 +\bigl[x=w_t\bigr]\bigl(\varepsilon(u_t)+g(u_t)\bigr),
\]
where $a(y)$ is the point that the pair of $y$ shares with $e_t$ and the sum
runs over both basis vectors of every pair in $F$.
\end{lemma}
\begin{proof}
No earlier pair contains both $i$ and $j$, so $F$ splits as $F_i\sqcup F_j$
and $V_{i,t}$ is spanned by the vectors of $F_i$, with
$v_{i,t}=\sum_{y\in F_i}\ell_i(y)$. The vectors with $g=1$ are linearly
independent, because each lies outside the span of all its predecessors,
and those belonging to pairs in $F$ lie in $V_{i,t}+V_{j,t}$. Hence
$\dim(V_{i,t}+V_{j,t})\ge\#\{y\in F:g(y)=1\}$ and
\[
 \dim(V_{i,t}\cap V_{j,t})
 =v_{i,t}+v_{j,t}-\dim(V_{i,t}+V_{j,t})
 \le\sum_{y\in F}\bigl(\ell_{a(y)}(y)-g(y)\bigr).
\]
This is the claim for $u_t$. For $w_t$ the spans gain $u_t$: the dimensions
of $V'_i,V'_j$ grow by $\ell_i(u_t),\ell_j(u_t)$, and the dimension of the
sum grows by at least $g(u_t)$, which adds $\varepsilon(u_t)+g(u_t)$.
\end{proof}

In particular a vector of type $(1,1,1)$ or $(0,0,0)$ creates no overlap
at any later pair; only defect vectors do, and only at pairs through the
point whose span they raised. In the direct case the lemma gives
$\pi(u_t)=0$ and $\pi(w_t)\le1$, recovering the exact exponent $-m$.

\begin{proposition}\label{prop:residual-28}
Within the source's nondirect residual stratum, for fixed even $D\ge28$,
the contribution is $O_D(q^{-D+28})$. In particular that stratum vanishes
for even $D\ge30$. The defect-one strata contribute $O_D(q^{-D-11})$.
All other dimension thresholds in the source remain required.
\end{proposition}
\begin{proof}
Fix the order of the active pairs obtained by listing the points as
$2,1,4,3,0,9,11,10,7,6,12,8,5$ and taking each pair at the later of its
two points, earlier points first. By Equation~\eqref{eq:sequential-remainder}
and Lemma~\ref{lem:overlap}, the remainder is at most
\begin{align*}
 R\le{}&-2m+\sum_{x\notin\mathcal D}\min\{N_i(x),N_j(x),T_i(x)+T_j(x)+[x=w_t]g(u_t)\}\\
 &+\sum_{x\in\mathcal D}\Bigl(N_i(x)+N_j(x)+[\text{type }(0,1,1)]\,\bigl(p(x)-\lceil d(x)/2\rceil\bigr)\Bigr),
\end{align*}
where $\mathcal D$ is the set of defect vectors, $N_a(x)$ counts the
earlier vectors at pairs through $a$, $T_a(x)=\sum\tau_a(y)$ over those
vectors, $p(x)$ is the position of $x$ in the list, and $d(x)$ is the
defect accumulated before $x$. Here $v'_a\le N_a$; the overlap of a
non-defect vector is at most the smaller endpoint span; the sum in
Lemma~\ref{lem:overlap} equals $T_i+T_j-\varepsilon(u_t)$ when $x=w_t$;
and $w'\le p-\lceil d/2\rceil$ because every defect vector lies in the
earlier global span and contributes at most two to the defect.

The right-hand side depends on the type sequence only through the counts
$T_a$ at points that still have unprocessed pairs, the previous vector's
$g$, and the defect so far. The chosen order keeps at most five points open
at any time, so an exact dynamic program over these quantities maximizes
the bound for every $d\le38$ through at most $2.2$ million states. The
maxima are $-11$ for $d=1$, $38$ for $d=2$, then $54$, $101$, $116$, $163$,
$177$, $222$, $235$ and $280$ for $d=3,\dots,10$; the ratio $R/d$ is
largest at $d=10$, where it equals $28$, and decreases afterwards. For
$d>38$ the averaging bound of Section~\ref{sec:sequential}, sharpened by
$w'\le2m-n_1-n_2$ for type $(0,1,1)$, gives
$R\le178+\tfrac{23}2n_1+(67-n_1-n_2)n_2$, whose ratio to $d$ is below $28$.
Hence $R\le28d$ for every $d\ge1$, and $-Dd+R\le-D+28$ once $D\ge28$. The
number of strata and type patterns is finite at fixed $D$, as before.
\end{proof}

The audit \texttt{audit\_overlap.py} implements this dynamic program in
standard-library Python: it reconstructs the pair order, exhausts all type
sequences for $d\le38$, and checks the tail bound. An integer program over
the same vector types, keeping the exact local-dimension accounting instead
of the relaxation $v'_a\le N_a$ and the exact count of earlier defect
vectors, returned the same maxima for every defect tried up to $26$ (not
part of the runner; the dynamic program is the certificate).

\paragraph{Reproducibility.}
The accompanying README documents \texttt{python -B run\_checks.py}.
The scripts, source manifest, and test expectations are in this repository.
The LaTeX source was compiled successfully with the Codex document editor.
\end{document}
